
\chapter{Methodology}
\label{ch:method}

This chapter presents the methodology adopted to investigate the research objectives
stated in Chapter~\ref{ch:into}. It describes the experimental design, the agent
architectures under comparison, the symlog transformation and its integration into
the learning pipeline, the environments used for evaluation, and the implementation
details necessary for reproducibility. The chapter concludes with a summary that
links the methodology to the results presented in Chapter~\ref{ch:results}.

% ================================================================
\section{Experimental Design}
\label{sec:method_design}
% ================================================================

The experiment isolates the contributions of two choices: (i) a learned world model for imagination-based training, and (ii) the application of the symlog transformation to critic and reward prediction targets. We define three configurations, each tested across two environments with five random seeds, yielding 30 runs. Table~\ref{tab:conditions} summarises the three conditions.

\begin{table}[!ht]
    \centering
    \caption{Summary of the three experimental conditions}
    \label{tab:conditions}
    \begin{tabular}{llccc}
        \toprule
        \textbf{Condition} & \textbf{Agent} & \textbf{World Model} & \textbf{Imagination} & \textbf{Symlog} \\
        \midrule
        A & Model-Free AC        & No  & No  & No  \\
        B & Imagination AC       & Yes & Yes & No  \\
        C & Imagination AC + Symlog & Yes & Yes & Yes \\
        \bottomrule
    \end{tabular}
\end{table}

Condition~A is the baseline model-free actor-critic agent. Condition~B adds a learned world model and uses imagined rollouts in a Dyna-style training loop. Condition~C retains this imagination pipeline and applies the symlog transformation to the critic value and world model reward predictions. Comparing A to B assesses the effect of imagination-based training; comparing B to C evaluates the effect of symlog.

We evaluate each configuration on CartPole-v1 and LunarLander-v3. We use five random seeds per configuration to account for variation in initialization and sampling. Performance is measured as the mean episodic return averaged over evaluation episodes.

% ================================================================
\section{Agent Architecture}
\label{sec:method_arch}
% ================================================================

All agents share a common neural network foundation built upon two-layer
multi-layer perceptrons (MLPs) with 256 hidden units per layer and ReLU
activations. A shared \texttt{build\_mlp} utility function constructs these
networks, ensuring architectural consistency across all components and conditions.
The three core components (actor, critic, and world model) are described below.

% ----------------------------------------------------------------
\subsection{Neural Network Components}
\label{subsec:method_nn}
% ----------------------------------------------------------------

Table~\ref{tab:nn_components} summarises the input, output, and training objective
of each neural network component.

\begin{table}[!ht]
    \centering
    \caption{Neural network components used across all conditions}
    \label{tab:nn_components}
    \begin{tabular}{lllll}
        \toprule
        \textbf{Component} & \textbf{Input} & \textbf{Output} & \textbf{Loss} & \textbf{Conditions} \\
        \midrule
        Actor   & State $s$ & Categorical $\pi(\cdot|s)$ & Policy gradient & A, B, C \\
        Critic  & State $s$ & Scalar $V(s)$              & MSE on TD target & A, B, C \\
        Dynamics head & $(s, a)$ & State residual $\Delta s$ & MSE & B, C \\
        Reward head   & $(s, a)$ & Scalar $\hat{r}$        & MSE & B, C \\
        \bottomrule
    \end{tabular}
\end{table}

The \textbf{actor} network maps a state vector to a categorical probability
distribution over the action space. During training, actions are sampled
from this distribution; during evaluation, the action with the highest probability
is selected. The actor is trained via policy gradient with an advantage baseline provided by the critic.

The \textbf{critic} network maps a state vector to a scalar estimate of the state
value $V(s)$. Under the standard configurations (A and B), the critic
outputs are in the original reward scale. Under the symlog configuration
(Condition~C), the critic outputs values in symlog space and the
\texttt{get\_value()} method applies the symexp inverse to recover
values in the original scale for advantage computation.

The \textbf{world model} contains two MLP heads that share no parameters. The
dynamics head takes a state-action pair and predicts a state residual $\Delta s$,
such that the predicted next state is $\hat{s}' = s + \Delta s$. This residual
formulation models only the change from the current state. The reward
head takes the same state-action pair and predicts the scalar reward. When symlog
is enabled (Condition~C), the reward head outputs predictions in symlog space and
the training target is transformed accordingly.

% ----------------------------------------------------------------
\subsection{Model-Free Actor-Critic (Condition~A)}
\label{subsec:method_mfac}
% ----------------------------------------------------------------

Condition~A implements a standard advantage actor-critic agent without a world
model. At each training step, a mini-batch of transitions
$(s, a, r, s', d)$ is sampled uniformly from the replay buffer. The one-step
temporal-difference (TD) target is computed as:
\begin{equation}
    \label{eq:td_target}
    y = r + \gamma (1 - d)\, V(s'),
\end{equation}
where $\gamma$ is the discount factor and $d \in \{0, 1\}$ indicates whether the
episode has terminated. The advantage is then:
\begin{equation}
    \label{eq:advantage}
    A(s, a) = y - V(s).
\end{equation}

The critic is updated by minimising the mean squared error between its prediction
and the TD target:
\begin{equation}
    \label{eq:critic_loss}
    \mathcal{L}_{\text{critic}} = \frac{1}{N} \sum_{i=1}^{N} \bigl( V(s_i) - y_i \bigr)^2.
\end{equation}

The actor is updated using the policy gradient objective with an entropy
regularisation term to encourage exploration:
\begin{equation}
    \label{eq:actor_loss}
    \mathcal{L}_{\text{actor}} = -\frac{1}{N} \sum_{i=1}^{N}
    \bigl[ \log \pi(a_i | s_i)\, A_i + \beta\, H(\pi(\cdot | s_i)) \bigr],
\end{equation}
where $H(\pi)$ denotes the entropy of the policy distribution and $\beta = 0.01$
is the entropy coefficient. The advantages $A_i$ are normalised to zero mean and
unit variance within each mini-batch to stabilise training.

% ----------------------------------------------------------------
\subsection{Imagination-Based Actor-Critic (Conditions~B and~C)}
\label{subsec:method_imag}
% ----------------------------------------------------------------

Conditions~B and~C extend the model-free baseline with a learned world model and
a Dyna-style training loop that supplements real experience
with imagined rollouts. The training procedure comprises two phases that execute at
each training step, plus an imagination phase that activates periodically.

\paragraph{Phase~1A: World Model Training.}
A mini-batch of real transitions is sampled from the replay buffer. The dynamics
head is trained to predict the state residual $\Delta s = s' - s$ via MSE loss,
and the reward head is trained to predict the observed reward. Two gradient steps
are performed per training call to accelerate world model convergence.

\paragraph{Phase~1B: Actor-Critic Training on Real Data.}
Using the same mini-batch, the actor and critic are updated identically to the
model-free procedure described in Section~\ref{subsec:method_mfac}. This ensures
that the agent continues to learn directly from real transitions regardless of the
quality of the world model.

\paragraph{Phase~2: Imagination Rollouts.}
Every 10th training step, after a warmup period of 1{,}000 environment steps, the
agent generates imagined experience. A batch of start states is sampled from the
replay buffer, and the world model is used to simulate $H = 5$ step rollouts from
each start state. At each imagined step, the actor selects an action from its
current policy, and the world model predicts the next state and reward. The
imagined trajectories are then used to compute $n$-step returns with critic
bootstrapping at the final step:
\begin{equation}
    \label{eq:nstep_return}
    G_t = \sum_{k=0}^{H-1} \gamma^k \hat{r}_{t+k} + \gamma^H V(\hat{s}_{t+H}),
\end{equation}
where $\hat{r}$ and $\hat{s}$ denote the world model predictions. The actor and
critic are then trained on this imagined data using the same loss functions as in
Phase~1B, with advantage normalisation applied independently to the imagined batch.

This approach operates directly in the raw observation
space rather than in a learned latent space, making it suitable for low-dimensional state spaces.

Algorithm~\ref{algo:imag_ac} presents the pseudocode for the complete training
loop of the imagination-based actor-critic.

\begin{algorithm}[!ht]
    \caption{Imagination-based actor-critic training loop}
    \label{algo:imag_ac}
    \begin{algorithmic}[1]
        \Require{Environment $\mathcal{E}$, replay buffer $\mathcal{B}$, networks $\pi_\theta$, $V_\phi$, $f_\psi$ (world model)}
        \Ensure{Trained policy $\pi_\theta$}
        \Statex
        \State Initialise $\mathcal{B}$, $\pi_\theta$, $V_\phi$, $f_\psi$ with random parameters
        \State $s \gets \mathcal{E}.\text{reset}()$
        \For{$t = 1$ to $T$}
            \State $a \sim \pi_\theta(\cdot | s)$ \Comment{Sample action from policy}
            \State $s', r, d \gets \mathcal{E}.\text{step}(a)$
            \State Store $(s, a, r, s', d)$ in $\mathcal{B}$
            \If{$|\mathcal{B}| \geq \text{batch\_size}$}
                \State Sample mini-batch $\{(s_i, a_i, r_i, s'_i, d_i)\}$ from $\mathcal{B}$
                \Statex \hspace{\algorithmicindent}\hspace{\algorithmicindent}\textbf{// Phase 1A: World model update (2 gradient steps)}
                \State Update $f_\psi$ on dynamics loss $\|\Delta \hat{s}_i - \Delta s_i\|^2$ and reward loss $\|\hat{r}_i - r_i\|^2$
                \Statex \hspace{\algorithmicindent}\hspace{\algorithmicindent}\textbf{// Phase 1B: AC update on real data}
                \State Compute TD targets $y_i = r_i + \gamma(1-d_i) V_\phi(s'_i)$
                \State Update $V_\phi$ via Eq.~\eqref{eq:critic_loss}, $\pi_\theta$ via Eq.~\eqref{eq:actor_loss}
                \Statex \hspace{\algorithmicindent}\hspace{\algorithmicindent}\textbf{// Phase 2: Imagination (every 10 steps, after warmup)}
                \If{$t \mod 10 = 0$ \textbf{and} $t > 1000$}
                    \State Sample start states $\{s^{(j)}_0\}$ from $\mathcal{B}$
                    \For{$k = 0$ to $H - 1$}
                        \State $a^{(j)}_k \sim \pi_\theta(\cdot | s^{(j)}_k)$
                        \State $\hat{s}^{(j)}_{k+1}, \hat{r}^{(j)}_k \gets f_\psi(s^{(j)}_k, a^{(j)}_k)$
                    \EndFor
                    \State Compute $n$-step returns $G^{(j)}$ via Eq.~\eqref{eq:nstep_return}
                    \State Update $V_\phi$, $\pi_\theta$ on imagined data
                \EndIf
            \EndIf
            \State $s \gets s'$ (reset if $d$)
        \EndFor
    \end{algorithmic}
\end{algorithm}

% ================================================================
\section{Symlog Integration}
\label{sec:method_symlog}
% ================================================================

The symlog transformation, introduced by \cite{hafner2023dreamerv3}, is a
symmetric logarithmic function designed to compress the scale of regression targets
while preserving their sign. It addresses a common challenge in reinforcement
learning: the reward scale varies dramatically across environments and over the
course of training, which can destabilise gradient-based optimisation when raw
values are used as regression targets.

\subsection{Mathematical Definition}
\label{subsec:method_symlog_def}

The symlog transformation and its inverse (symexp) are defined as:
\begin{equation}
    \label{eq:symlog_meth}
    \text{symlog}(x) = \text{sign}(x) \cdot \ln(|x| + 1),
\end{equation}
\begin{equation}
    \label{eq:symexp_meth}
    \text{symexp}(x) = \text{sign}(x) \cdot \bigl( e^{|x|} - 1 \bigr).
\end{equation}
These functions satisfy $\text{symexp}(\text{symlog}(x)) = x$ for all $x$. The
symlog function behaves approximately linearly near zero and logarithmically for
large magnitudes, providing a smooth compression that is well-suited to targets
whose scale is not known in advance.

\subsection{Application Points}
\label{subsec:method_symlog_apply}

In Condition~C, the symlog transformation is applied at two points in the learning
pipeline. Table~\ref{tab:symlog_application} contrasts the standard and symlog
configurations.

\begin{table}[!ht]
    \centering
    \caption{Application of the symlog transformation across conditions}
    \label{tab:symlog_application}
    \begin{tabular}{lll}
        \toprule
        \textbf{Component} & \textbf{Standard (A, B)} & \textbf{Symlog (C)} \\
        \midrule
        Critic output space   & Original scale       & Symlog space \\
        Critic TD target      & $y = r + \gamma V(s')$ & $\text{symlog}(r + \gamma\, \text{symexp}(V_{\text{raw}}(s')))$ \\
        Critic \texttt{get\_value()} & $V(s)$        & $\text{symexp}(V_{\text{raw}}(s))$ \\
        WM reward target      & $r$                  & $\text{symlog}(r)$ \\
        WM reward prediction  & $\hat{r}$            & $\hat{r}_{\text{raw}}$ (in symlog space) \\
        WM reward at inference & $\hat{r}$           & $\text{symexp}(\hat{r}_{\text{raw}})$ \\
        \bottomrule
    \end{tabular}
\end{table}

When symlog is enabled, the critic network learns to predict values in the
compressed symlog space. The MSE loss is computed between the raw network output
and the symlog-transformed TD target. At inference time, the \texttt{get\_value()}
method applies symexp to recover the value estimate in the original reward scale,
which is then used for advantage computation. Similarly, the world model reward
head predicts rewards in symlog space; during imagination rollouts, symexp is
applied to convert predicted rewards back to the original scale before computing
$n$-step returns.

\subsection{Numerical Stability}
\label{subsec:method_stability}

We use two measures to ensure numerical stability when using the symlog and
symexp transformations. First, the input to symexp is clamped to the range
$[-20, 20]$ to prevent overflow in the exponential computation.
Second, gradient norms are clipped to a
maximum of 1.0 for all networks across all conditions. This gradient
clipping serves as a safeguard against training instability from
large value predictions, particularly during early training when the world model
is still inaccurate.

% ================================================================
\section{Environment Details}
\label{sec:method_envs}
% ================================================================

We evaluate the agents on two Gymnasium environments that differ in dimensionality,
action space size, and reward structure: CartPole-v1 and LunarLander-v3.

\subsection{CartPole-v1}
\label{subsec:method_cartpole}

CartPole-v1 is a classic control task in which the agent must balance a pole on a
cart by applying left or right forces. The observation space is a four-dimensional
continuous vector comprising cart position, cart velocity, pole angle, and pole
angular velocity. The action space is discrete with two actions (push left, push
right). A reward of $+1$ is received at every time step the pole remains upright,
and the episode terminates when the pole angle exceeds $\pm 12^\circ$, the cart
moves beyond $\pm 2.4$ units from the centre, or 500 steps are reached. The
maximum achievable return is therefore 500. CartPole-v1 serves as a simple
baseline environment with a narrow, non-negative reward range, where symlog
compression is expected to have minimal impact.

\subsection{LunarLander-v3}
\label{subsec:method_lunarlander}

LunarLander-v3 is a more challenging control task in which the agent must land a
spacecraft on a designated pad. The observation space is an eight-dimensional
continuous vector encoding position, velocity, angle, angular velocity, and leg
contact indicators. The action space is discrete with four actions (do nothing,
fire left engine, fire main engine, fire right engine). The reward structure is
complex: the agent receives positive reward for moving toward the landing pad and
for leg contact, but incurs penalties for fuel usage and crashes. Typical episodic
returns range from approximately $-200$ (crash) to $+300$ (successful landing).
LunarLander-v3 is included because its wide, mixed-sign reward range creates the
conditions under which symlog compression is expected to be most beneficial.

\subsection{Environment Selection Rationale}
\label{subsec:method_env_rationale}

CartPole-v1 provides a low-dimensional, narrow-reward setting to verify basic functionality. LunarLander-v3 features higher dimensionality, a larger action space, and a variable reward signal to evaluate the benefits of imagination-based training and symlog scaling.

% ================================================================
\section{Implementation Details}
\label{sec:method_impl}
% ================================================================

\subsection{Software Framework}
\label{subsec:method_framework}

All agents are implemented in Python using PyTorch for
neural network construction, automatic differentiation, and optimisation. The
environments are instantiated via the Gymnasium API.
Training was conducted on CPU, as the network architectures and environments are computationally lightweight.

\subsection{Hyperparameters}
\label{subsec:method_hyperparams}

Table~\ref{tab:hyperparams} lists the hyperparameters shared across all conditions.
Where a hyperparameter applies only to conditions with a world model (B and~C),
this is indicated in the table.

\begin{table}[!ht]
    \centering
    \caption{Hyperparameters used across all experimental conditions}
    \label{tab:hyperparams}
    \begin{tabular}{llll}
        \toprule
        \textbf{Hyperparameter} & \textbf{Value} & \textbf{Conditions} \\
        \midrule
        Learning rate (all networks)        & $3 \times 10^{-4}$   & A, B, C \\
        Discount factor $\gamma$            & 0.99                 & A, B, C \\
        Batch size                          & 64                   & A, B, C \\
        Replay buffer capacity              & 100{,}000            & A, B, C \\
        Entropy coefficient $\beta$         & 0.01                 & A, B, C \\
        Gradient clip norm                  & 1.0                  & A, B, C \\
        Hidden layer size                   & 256                  & A, B, C \\
        Number of hidden layers             & 2                    & A, B, C \\
        Activation function                 & ReLU                 & A, B, C \\
        Imagination horizon $H$             & 5                    & B, C \\
        Warmup steps before imagination     & 1{,}000              & B, C \\
        Imagination frequency               & Every 10 steps       & B, C \\
        WM gradient steps per update        & 2                    & B, C \\
        Symexp clamp range                  & $[-20, 20]$          & C \\
        \midrule
        Training steps (CartPole-v1)        & 100{,}000            & A, B, C \\
        Training steps (LunarLander-v3)     & 300{,}000            & A, B, C \\
        \bottomrule
    \end{tabular}
\end{table}

All networks use the Adam optimiser with the same learning rate. No learning rate
scheduling or decay is applied. The hyperparameters were selected based on
defaults from the reinforcement learning
literature and kept constant across all conditions to ensure a
fair comparison.

\subsection{Replay Buffer}
\label{subsec:method_buffer}

A fixed-capacity replay buffer stores transitions as pre-allocated NumPy arrays.
When the buffer is full, new transitions overwrite the oldest entries in a
circular fashion. During training, mini-batches are sampled uniformly at random.
The sampled arrays are converted to PyTorch tensors via zero-copy
\texttt{torch.from\_numpy()} calls, avoiding memory allocation and
improving sampling throughput.

\subsection{Evaluation Protocol}
\label{subsec:method_eval}

Agent performance is evaluated every 1{,}000 environment steps by running 10
evaluation episodes with the greedy policy (i.e., the action with the highest
probability under the actor is selected deterministically). The mean and standard
deviation of the episodic returns across these 10 episodes are recorded. This
protocol is identical across all configurations. At reporting time, results are aggregated across
the five seeds per configuration.

% ================================================================
\section{Summary}
\label{sec:method_summary}
% ================================================================

This chapter described the methodology for evaluating imagination-based training and symlog-transformed targets. We defined three configurations (A, B, and C) sharing the same network architecture, hyperparameters, and evaluation protocol. Chapter~\ref{ch:results} presents the experimental results.